% Ember -- full build plan, FIRE-SPREAD version (Qiskit Fall Fest Waterloo 2026, team of 5, present Fri 6 Nov 2026).
% Replaces the FWI-forecasting plan (kept as ember-plan-fwi.tex = fallback if the perimeter check returns verdict C).
% Style follows project3_QRC_forecasting.pdf: every load-bearing claim carries a provenance tag.
% Sources: CWFIS fwi_obs files + station extract (4 Nov 2026); CWFIS hotspots/archive listing (perimeters 2018--2023);
%   PronoiaQ code (~/Projects/Quantathon2026-QRC_detection); github.com/Luke-Fu101/QRC_Broadcast_Burns (read 4 Nov 2026);
%   hackathon notes pack (chapters 2, 3, 7, 8); docs/perimeter-check.md (verdict B, re-checked 4 Nov 2026).
% Single-pass safe: no \ref, no \tableofcontents. Contents page links to hyperref's own anchors (section.N).
%   Deliver .tex only; do not compile here.

\documentclass[11pt,a4paper]{article}
\usepackage[margin=2.2cm,headheight=14pt]{geometry}
\usepackage[T1]{fontenc}
\usepackage{lmodern}
\usepackage{microtype}
\usepackage{amsmath,amssymb}
\usepackage{booktabs,array,tabularx}
\usepackage[dvipsnames]{xcolor}
\usepackage[most]{tcolorbox}
\usepackage{enumitem}
\usepackage{fancyhdr}
\usepackage{titlesec}
\usepackage[colorlinks=true,linkcolor=NavyBlue,urlcolor=NavyBlue]{hyperref}
\usepackage{bookmark}

\setlist{nosep,leftmargin=*,topsep=2pt}
\setlength{\parindent}{0pt}
\setlength{\parskip}{5pt}
\renewcommand{\arraystretch}{1.3}
\newcolumntype{Y}{>{\raggedright\arraybackslash}X}

% ---- section look ----
\titleformat{\section}{\Large\bfseries\color{NavyBlue!85!black}}{\thesection}{0.6em}{}[\vspace{-2pt}{\color{NavyBlue!40}\titlerule}]
\titleformat{\subsection}{\large\bfseries}{\thesubsection}{0.5em}{}
\titlespacing*{\section}{0pt}{14pt}{6pt}
\titlespacing*{\subsection}{0pt}{10pt}{4pt}

% ---- header / footer ----
\pagestyle{fancy}
\fancyhf{}
\fancyhead[L]{\small\nouppercase{\leftmark}}
\fancyhead[R]{\small Ember -- build plan (fire spread)}
\fancyfoot[C]{\small\thepage}
\renewcommand{\headrulewidth}{0.4pt}
\renewcommand{\sectionmark}[1]{\markboth{\thesection\quad #1}{}}

% ---- boxes ----
\tcbset{pb/.style={enhanced,breakable,arc=2pt,boxrule=0.6pt,left=6pt,right=6pt,top=4pt,bottom=4pt,fonttitle=\bfseries}}
\newtcolorbox{glance}{pb,colback=NavyBlue!4,colframe=NavyBlue!80!black,title={Plan at a glance}}
\newtcolorbox{watch}{pb,colback=red!4,colframe=red!65!black,title={Watch out}}
\newtcolorbox{decision}{pb,colback=Goldenrod!10,colframe=Goldenrod!70!black,title={Design decision}}
\newtcolorbox{teamrule}{pb,colback=ForestGreen!5,colframe=ForestGreen!70!black,title={Team rule}}
\newtcolorbox{gate}{pb,colback=Orange!6,colframe=Orange!80!black,title={Go / no-go gate (Wed night)}}

% ---- provenance tags (short) ----
\newcommand{\tagbox}[2]{{\footnotesize\color{#1}\textsf{[#2]}}}
\newcommand{\supp}[1]{\tagbox{ForestGreen!70!black}{S: #1}}  % supported by a source / earlier run
\newcommand{\veri}[1]{\tagbox{ForestGreen!70!black}{V: #1}}  % verified today
\newcommand{\ours}{\tagbox{NavyBlue}{D}}                      % our design choice
\newcommand{\caut}[1]{\tagbox{BrickRed}{!: #1}}               % caution
\newcommand{\todo}[1]{\tagbox{Orange!85!black}{?: #1}}        % owner must check
\newcommand{\code}[1]{\texttt{\small #1}}
\newcommand{\tocsec}[2]{\par\noindent\hyperlink{section.#1}{\textbf{#1}\quad #2}\par}

\begin{document}

% ================= page 1 =================
\begin{center}
{\LARGE\bfseries Ember}\par\smallskip
{\Large Quantum reservoir forecasting of wildfire spread}\par\medskip
{\normalsize Full build plan \quad$\cdot$\quad Qiskit Fall Fest Waterloo 2026}\par
{\small Team of 5 \quad$\cdot$\quad build Wed 4 -- Thu 5 Nov \quad$\cdot$\quad present Fri 6 Nov, Davis Centre}\par
\end{center}

\begin{glance}
\begin{tabularx}{\linewidth}{@{}l Y@{}}
\textbf{Goal} & From day 3 of an active fire, predict how far the fire edge moves in each of 8 compass sectors over the
next 24\,h and 72\,h. Turn that into a predicted outline and a time-to-reach for nearby towns.\\
\textbf{Data} & CWFIS satellite hotspots 2018--2023. Daily fire outlines are \textbf{rebuilt} from the hotspots;
each hotspot also carries its local weather and fire-danger codes.\\
\textbf{Model} & Fixed 5-qubit Qiskit circuit (reservoir) $\to$ measured features $\to$ ridge read-out. Only the
read-out is trained (42 numbers).\\
\textbf{Compare} & Persistence, physics spread rate (FBP), ridge on raw window, same-size classical reservoir (ESN),
random circuit.\\
\textbf{Hardware} & One batched IBM job: fire ZF015 (Yellowknife), 10--20 Aug 2023, 11 days $\times$ 8 sectors = 88
circuits.\\
\textbf{Result} & Reservoir $\approx$ classical reservoirs, slightly worse. FBP lowest mean error; persistence best outline overlap; learned models best on growing sectors. \textbf{No quantum advantage.}\\
\textbf{Credit} & Extends PronoiaQ's reservoir (Quantathon 2026, 2nd place + Best Team).\\
\end{tabularx}
\end{glance}

\begin{gate}
\textbf{Result: verdict B.} \veri{docs/perimeter-check.md; key numbers re-run independently}
\begin{itemize}
\item The perimeter archives hold \textbf{one final outline per fire per season} (2023: 1{,}144 outlines, 1 date
  each). No daily history. Fires with $\ge10$ dated outlines: 0 in every year.
\item The hotspot archive \textbf{is} dated: every detection has a UTC timestamp. Fire ZF015 (Yellowknife) has hotspots
  on 70 of 77 days, all 11 days of 10--20 Aug.
\item $\Rightarrow$ R1 rebuilds daily outlines from hotspots (Section 3). Target = reconstructed advance, not an
  independently observed edge. Say so on the slide.
\item If reconstruction fails by Thu 09:00: switch to \code{ember-plan-fwi.tex}. It is ready to run.
\end{itemize}
\end{gate}

\subsection*{Contents}
\tocsec{1}{Problem and value}
\tocsec{2}{What is reused, what is new}
\tocsec{3}{Data}
\tocsec{4}{Method}
\tocsec{5}{Baselines and controls}
\tocsec{6}{Experiments}
\tocsec{7}{Hardware plan}
\tocsec{8}{Team, interfaces, timeline}
\tocsec{9}{Expected outcomes}
\tocsec{10}{Results (test 2023, scored once)}
\tocsec{11}{Pitch and questions}
\tocsec{12}{Soundness audit}
\par\noindent\hyperlink{refs}{\textbf{--}\quad References}\par

\subsection*{How to read the tags}
\begin{center}\small
\begin{tabular}{@{}ll@{}}
\supp{source} & backed by a source or our own earlier run\\
\veri{what} & checked on 4 Nov 2026\\
\ours & our design choice\\
\caut{why} & risk, or not supported\\
\todo{who} & owner must confirm\\
\end{tabular}
\end{center}

\newpage
% =====================================================================
\section{Problem and value}

\begin{itemize}
\item 2023 season: $\approx17.2$\,M ha burned. \supp{CWFIS; notes ch.~1}
\item Aug 2023: all of Yellowknife evacuated ahead of a fire. \caut{cite a primary source on the slide}
\item Once a fire is burning, the question is \textbf{where it goes next}: which side grows, how fast, and which town
  it reaches first.
\item Evacuations are planned in 24--72\,h steps: planes, roads, crews. \caut{general practice; find a source}
\item Under-predicting spread = late evacuation. $\Rightarrow$ report errors on the fast-growing side separately. \ours
\end{itemize}

\textbf{Question we answer:}
\begin{enumerate}
\item Given the last 5 days of a fire's growth and local weather, how far will its edge move in each compass
  sector in the next 24\,h and 72\,h?
\item Does a fixed quantum reservoir extract that signal as well as a classical model of the same size?
\end{enumerate}

\textbf{Not in scope:} where fires start (ignition), fine-grained spread maps (operational fire-spread models already do
this from fuel maps), and evacuation routing.

% =====================================================================
\section{What is reused, what is new}

\begin{center}\small
\begin{tabularx}{\linewidth}{@{}Y l l@{}}
\toprule
\textbf{Piece} & \textbf{Status} & \textbf{Where}\\
\midrule
Windowed restart reservoir (encode, entangle, frozen rotations) & reused & \code{QRC\_theoretical.py}\\
Qiskit twin, endianness fix, Aer noise path & reused & \code{QRC\_QPU\_implementation.py}\\
One ridge read-out per lead time & reused & \code{QRC\_horizon\_forecast/}\\
CWFIS station extract (fallback weather, FWI plan) & reused & \code{scripts/extract\_stations.sh}\\
Perimeter / hotspot check & reused & \code{scripts/check\_perimeters.py}\\
Daily outlines rebuilt from hotspots; 8-sector edge geometry & \textbf{new} & --\\
Wind-along-sector input & \textbf{new} & --\\
Predicted outline + time-to-reach for towns & \textbf{new} & --\\
Chain entangler (no SWAPs on IBM chips) & \textbf{new} & --\\
Real IBM hardware run & \textbf{new} & PronoiaQ never ran it\\
\bottomrule
\end{tabularx}
\end{center}

\begin{watch}
PronoiaQ was a team of four. Say ``extends PronoiaQ's reservoir''. Do not present the engine as this week's work.
\end{watch}

\textbf{Related work to name on the slide.} \code{QRC\_Broadcast\_Burns} (GitHub, Aug 2026) tests a reservoir for
prescribed-burn weather windows, but only on the Lorenz-63 toy system, simulated exactly, with no fire data.
\veri{repo read 4 Nov} Ember uses real fire perimeters, a physical target in km, and a hardware run. \ours

% =====================================================================
\section{Data}

\subsection{Hotspots (the main data)}
\begin{itemize}
\item Archive: \url{https://cwfis.cfs.nrcan.gc.ca/downloads/hotspots/archive/}, \code{<YEAR>\_hotspots.zip}
  (2012--2023, all HTTP 200). 2023: one CSV, 3{,}674{,}409 rows, 35 columns, 748\,MB unzipped. \veri{perimeter-check.md}
\item 2018--2022 zips downloaded into \code{data/raw/cwfis\_perimeters/} (sizes match the server). Same columns
  every year; no negative weather values in Canada. \veri{U1b}
\item Each row = one satellite detection (MODIS or VIIRS). \code{rep\_date} = UTC time. \code{lat, lon} in degrees.
  \veri{CSV}
\item Each row also carries local-noon weather and codes: \code{ws} (km/h), \code{wd}, \code{rh}, \code{isi},
  \code{bui}, plus modelled \code{ros} (spread rate, m/min) and \code{fuel} type.
  \supp{NRCan hotspot attribute sheet}
\item Slides must credit: Canadian Forest Service, CWFIS. \supp{CWFIS metadata}
\end{itemize}

\subsection{Final perimeters (grouping only)}
\begin{itemize}
\item \code{<YEAR>\_perimeters.zip}, 2018--2023: one polygon per fire for the whole season, fields
  \code{uid, hcount, firstdate, lastdate, area} (+\code{consis\_id} in 2023). \veri{all six files}
\item \code{firstdate}/\code{lastdate} = first and last contributing hotspot. They are \textbf{not} perimeter dates.
\item Use: assign each hotspot to a fire (point inside the final polygon). Nothing else.
\item ZF015-23 (Yellowknife) = \code{uid} 242981: 29 Jun -- 13 Sep, 1{,}781\,km$^2$ (geometry), final outline
  24.9\,km from the town. \veri{re-run}
\end{itemize}

\begin{watch}
Grouping by the final polygon uses the end-of-season footprint, so it knows which hotspots belong together. It never
uses their timing: the outline on day $t$ is built only from detections up to day $t$. State this rule on the
method slide.
\end{watch}

\subsection{Rebuilding the daily outline}
For fire $f$, let $H_f(t)$ = all its hotspots detected up to the end of UTC day $t$.
\begin{enumerate}
\item Drop isolated detections: points with no other detection of $f$ within 2\,km. \ours
\item Seed = largest cluster of day-1 hotspots (single linkage, 2\,km). Centre $c$ = its mean position. Fixed for the
  life of the fire. \ours
\item A hotspot joins the fire only through a chain of hotspots each $\le2$\,km apart, starting from the seed. \ours
\item After day 1, a joining hotspot counts only if it lies within yesterday's edge $+30$\,km in its wedge;
  otherwise it never counts (another fire merging in). \ours
\item Wedge $k$ = bearings $45^\circ k\pm22.5^\circ$ from $c$.
\end{enumerate}
\[
r_k(t)=\max_{p\,\in\,H_f(t),\ p\,\in\,\text{wedge }k}\ \lVert p-c\rVert\ \ \text{(km)},
\qquad
\Delta d_k(t)=r_k(t)-r_k(t-1)\ \ge 0
\]
Because $H_f(t)$ only grows, $r_k$ never shrinks (burned ground stays burned). The 8 points $r_k$ joined around $c$
form the outline used for the map and for IoU. \ours

\begin{center}\small
\begin{tabularx}{\linewidth}{@{}l Y@{}}
\toprule
\textbf{Limitation} & \textbf{Effect / what we do}\\
\midrule
Cloud, smoke, satellite gaps & days with no detection $\Rightarrow$ $\Delta d=0$ that day, jump later. Keep a
missing-day mask; E9 tests this.\\
Pixel size 375\,m (VIIRS) -- 1\,km (MODIS) & edge error $\sim$0.5\,km. Report MAE in km next to this.\\
Inputs and targets from the same hotspots & skill is against reconstructed edges, not ground truth. Say so.\\
Several ignitions in one perimeter & without rules 3--4, one-day jumps up to 144\,km. With them: max 30\,km/day;
2023: 4.3\% of hotspots never count, in 23 of 392 fires. A merged fire can still creep in 30\,km steps. \veri{U1c}\\
Rebuilt vs official final outline & E10, 2023: median IoU 0.65 (8-point outline), 0.83 (buffered hotspots).
\veri{e10\_2023.json}\\
\bottomrule
\end{tabularx}
\end{center}

\subsection{Weather at the fire}
\begin{itemize}
\item Use the hotspot rows themselves: on day $t$, the \textbf{median} \code{ws, wd, rh, isi, bui} over that day's
  hotspots of fire $f$. No station distance problem. \ours
\item No hotspot that day $\Rightarrow$ carry the last value, flag the day.
  \caut{weather missing on 23 of 3.67\,M rows; rows with RH $=0$ are placeholders, drop them}
\item Fallback: nearest federal station from \code{fwi\_obs/} (already downloaded). \veri{station extract}
\item \code{wd} = direction the wind blows \emph{from}, in degrees. \supp{standard meteorological convention}
\end{itemize}

\subsection{Wind along each sector}
Sector $k$ points at bearing $\varphi_k=45^\circ k$ (N, NE, \dots, NW). Wind component pushing toward sector $k$:
\[
w_k(t)=\mathrm{ws}(t)\,\cos\!\big(\varphi_k-(\mathrm{wd}(t)+180^\circ)\big)
\]
Positive = wind blows the fire toward that sector.

\begin{decision}
\textbf{8 sectors, not a pixel map.} Hotspot outlines are coarse (0.4--1\,km pixels, daily), and fuel-map spread
models already produce pixel maps. Eight numbers per fire-day keep the target simple, give $8\times$ the samples, and fit
a 5-qubit read-out. Fix the centre on day 1 so $r_k$ measures real growth, not a moving origin. \ours
\end{decision}

\subsection{Samples, splits, targets}
\begin{center}\small
\begin{tabularx}{\linewidth}{@{}l Y@{}}
\toprule
\textbf{Item} & \textbf{Rule}\\
\midrule
Sample & one fire $\times$ day $t$ $\times$ sector $k$ \ours\\
Fires kept & $\ge10$ hotspot days and final area $\ge10$\,km$^2$; no merge with another uid \ours\\
Window & days $t-4,\dots,t$ ($L=5$), oldest first \ours\\
Before day 1 & $\Delta d_k=0$ (fire did not exist), weather as observed. Forecasts start on day 3. \ours\\
Targets & cumulative advance $y^{(1)}=r_k(t+1)-r_k(t)$ and $y^{(3)}=r_k(t+3)-r_k(t)$, in km \ours\\
Train & fires that started in 2018--2021\\
Validation & 2022\\
Test & 2023 (includes the Yellowknife-area fire)\\
Split rule & By year, so one fire never sits in two splits. Never shuffle fire-days. \supp{notes ch.~3}\\
Scaling & Min--max on training rows, clip to $[0,1]$. \ours\\
Gaps & Day without hotspots: $\Delta d_k=0$, weather carried, mask set. Drop windows with $>2$ masked days. \ours\\
\bottomrule
\end{tabularx}
\end{center}

\begin{watch}
\begin{itemize}
\item The 8 sectors of one fire-day are not independent, and neither are consecutive days of one fire.
  Bootstrap \textbf{by fire}, not by row, for every error bar.
\item 2023 was an extreme year. Test errors may be larger than validation errors. Report both; do not retune on test.
\item Usable fires: train 378 (2018: 134, 2019: 44, 2020: 7, 2021: 193), val 102, test 392. Forecast fire-days:
  14{,}654 / 4{,}763 / 20{,}248. \veri{fire\_counts.json}
\end{itemize}
\end{watch}

% =====================================================================
\section{Method}

\subsection{Circuit (one per sample, restarted each time)}
Feed the window $\mathbf u_{t-4},\dots,\mathbf u_t$, oldest first. Each step $j$ applies:
\[
\underbrace{\bigotimes_i R_Y\!\big(\gamma\, g_i\, u_{j,i}\big)}_{\text{encode}}
\;\longrightarrow\;
\underbrace{\prod_{i=0}^{n-2} \mathrm{CZ}_{i,i+1}}_{\text{chain entangler}}
\;\longrightarrow\;
\underbrace{\bigotimes_i R_Z(b_i)\,R_Y(a_i)}_{\text{frozen random}}
\]
The encoding rotation adds to the state left by earlier days; it does not overwrite it. After day $t$ the state holds
the whole window, recent days strongest. Measuring ends the run; the next sample starts again from $|00000\rangle$.

\begin{center}\small
\begin{tabularx}{\linewidth}{@{}l l Y@{}}
\toprule
\textbf{Symbol} & \textbf{Value} & \textbf{Meaning}\\
\midrule
$n$ & 5 & qubits, one per input \supp{PronoiaQ}\\
$L$ & 5 & window in days; E3 swept 3, 5, 7 on 2022, chose 5 \veri{frozen\_config.json}\\
$\gamma$ & $\pi/4$ & encoding gain; E3 swept $\pi/4,\pi/2,3\pi/4$ on 2022, chose $\pi/4$ \veri{frozen\_config.json}\\
$g_i,a_i,b_i$ & random, seed 7 & frozen, never trained \supp{PronoiaQ}\\
\bottomrule
\end{tabularx}
\end{center}

\subsection{Input on each qubit}
\begin{center}\small
\begin{tabularx}{\linewidth}{@{}l l Y@{}}
\toprule
\textbf{Qubit} & \textbf{Input} & \textbf{Why here}\\
\midrule
q0 & wind along sector $w_k$ & fastest driver; pushes the head of the fire\\
q1 & ISI (initial spread index) & wind $\times$ fine-fuel dryness; next to wind \supp{Van Wagner 1987}\\
q2 & recent advance $\Delta d_k$ & the fire's own state; middle of the chain couples to all drivers\\
q3 & BUI (buildup index) & fuel available to burn; changes over weeks \supp{Van Wagner 1987}\\
q4 & RH (relative humidity) & drying power of the air; never missing\\
\bottomrule
\end{tabularx}
\end{center}
\ours\ Sweep one alternative map in E3 (q0 $w_k$, q1 $\Delta d_k$, q2 ISI) so the order is chosen, not assumed.

\begin{decision}
\textbf{Chain, not ring.} PronoiaQ entangles a 5-qubit ring. IBM heavy-hex chips have no 5-cycles, so a ring
needs SWAPs. A chain maps to a path on the chip with \textbf{zero SWAPs}. Use the chain in simulation too, so
simulation and hardware are the same model. \ours
\end{decision}

\subsection{Features}
\begin{center}\small
\begin{tabular}{@{}lll@{}}
\toprule
\textbf{Basis} & \textbf{Features} & \textbf{Count}\\
\midrule
Z & $\langle Z_i\rangle$, \ $\langle Z_iZ_j\rangle$ & $5+10=15$\\
X (extra circuit) & $\langle X_i\rangle$ & 5\\
\midrule
Simulation & Z + X + bias & 21\\
Hardware & Z + bias & 16\\
\bottomrule
\end{tabular}
\end{center}
\supp{PronoiaQ \code{counts\_to\_features}}

\begin{watch}
Qiskit count keys are little-endian (qubit 0 = rightmost bit). Reverse exactly once, inside
\code{counts\_to\_features}. Nowhere else.
\end{watch}

\subsection{Read-out (the only trained part)}
\[
w_h^\star=(X^\top X+\lambda I)^{-1}X^\top y^{(h)},
\qquad
\hat y^{(h)}=w_h^{\star\top}\,x,
\qquad h\in\{1,3\}
\]
\begin{itemize}
\item One ridge per lead: $2\times21=42$ trained numbers in total.
\item $\lambda$ from validation, grid $10^{-6}$ to $10^{2}$.
\item No gradients through the circuit $\Rightarrow$ no barren plateaus, training takes seconds. \supp{notes ch.~3}
\end{itemize}

\begin{center}\small
\begin{tabularx}{\linewidth}{@{}l l Y@{}}
\toprule
\textbf{Part} & \textbf{Trained?} & \textbf{How it is set}\\
\midrule
Fire tracks, sectors, weather join & no & geometry and table joins\\
Scaling to $[0,1]$ & fitted & min/max of training rows only\\
Reservoir gates & \textbf{never} & random, seed 7, frozen\\
$\gamma$, $L$, input map, $\lambda$ & tuned & on 2022 validation, then frozen\\
\textbf{Ridge weights} $w_1,w_3$ & \textbf{yes} & least squares on 2018--2021\\
Outline and time-to-reach & no & arithmetic on predictions\\
\bottomrule
\end{tabularx}
\end{center}

\subsection{From predictions to a map and a time-to-reach}
\[
\hat r_k(t+h)=r_k(t)+\hat y_k^{(h)},
\qquad
\widehat{\mathrm{ETA}}_{\text{town}}\approx\frac{D-r_k(t)}{\hat y_k^{(1)}}\ \text{days}
\]
Join the 8 predicted points into the predicted outline. $D$ = distance from the centroid to the town, $k$ = the sector
the town lies in. Show the time-to-reach only when $\hat y_k^{(1)}>0$. \ours

\caut{ZF015 (E8): the single-sector formula gives 80--1355 days for Yellowknife (66.7\,km, sector 3). The fire's
closest approach is through the south lobe (sector 4), about 25\,km away. Do not show a single-sector ETA; show the
outline and the closest-approach distance instead.}

% =====================================================================
\section{Baselines and controls}

\begin{teamrule}
Build and score the baselines \textbf{before} the quantum reservoir.
\end{teamrule}

\begin{center}\small
\begin{tabularx}{\linewidth}{@{}l Y Y@{}}
\toprule
\textbf{Model} & \textbf{Definition} & \textbf{Why it is here}\\
\midrule
Persistence & $\hat y^{(h)}=h\,\Delta d_k(t)$ (keep today's speed) & the bar to beat\\
Mean advance & mean of $y^{(h)}$ in training rows & reference for skill\\
Physics (FBP) & $\hat y^{(h)}_k=h\cdot\mathrm{ros}\cdot T_b\cdot\max(0,\cos\theta_k)$; \ $\mathrm{ros}$ = median modelled
spread rate of that day's hotspots, $\theta_k$ = angle between sector and downwind, $T_b$ = burning minutes per day,
fitted on training & what the operational model says \supp{hotspot \code{ros} field} \caut{crude: head-fire rate only}\\
Raw window & ridge on the raw 5-day inputs (25 numbers) & does the reservoir add anything? \ours\\
ESN & tanh reservoir, 5 and 20 units, \textbf{fully connected}, same ridge & same read-out size \supp{notes ch.~8}\\
Random circuit & same circuit, Haar-random entangler, 10 seeds & tuned dynamics vs just ``quantum''
\supp{notes ch.~8}\\
Reset reservoir & Fujii--Nakajima style: reset input qubits each day, hidden qubits keep memory (simulation only) &
the other standard reservoir design \supp{Fujii--Nakajima 2017} \caut{not run: time}\\
\bottomrule
\end{tabularx}
\end{center}

\begin{watch}
\code{reservoirpy} sets connectivity to 0.1 by default. With 5 units most units then have no inputs, and the ESN loses
for a silly reason. Set input and recurrent connectivity to 1.0. \supp{QRC\_Broadcast\_Burns README}
\end{watch}

Scores per lead $h$:
\[
\mathrm{MAE}_h=\frac1N\sum\big|y^{(h)}-\hat y^{(h)}\big|\ \text{(km)},
\qquad
\mathrm{SS}_h=1-\frac{\mathrm{MAE}_h^{\text{model}}}{\mathrm{MAE}_h^{\text{persist}}},
\qquad
\mathrm{IoU}=\frac{|\hat P\cap P|}{|\hat P\cup P|}
\]
$\mathrm{SS}_h\le0$ means no better than persistence. $P,\hat P$ = observed and predicted outlines. Also report
$\mathrm{MAE}$ on growing rows, $y>0$ (where under-prediction is dangerous). 91\% of rows do not grow, so a ``fastest quarter'' would be all zeros. \ours

\begin{watch}
\textbf{Operational spread models know more than we do.} Canadian fire centres run the Fire Behaviour Prediction
system and spread simulators with fuel maps and weather forecasts. Ember has no fuel map and no forecast. The FBP
baseline row shows where we stand against the physics. Position Ember as ``a data-driven check from satellite
detections only''. Rehearse this answer.
\end{watch}

% =====================================================================
\section{Experiments}

\begin{center}\small
\begin{tabularx}{\linewidth}{@{}l Y Y l@{}}
\toprule
\textbf{ID} & \textbf{Run} & \textbf{Output} & \textbf{Owner}\\
\midrule
E1 & All models, exact features, test 2023, both leads & main table: MAE, SS, IoU, fast-quarter MAE & R3\\
E2 & Error by sector relative to wind (downwind vs upwind) & ``where it works'' plot & R5\\
E3 & Sweep $L$, $\gamma$, input map on validation & frozen config & R2\\
E4 & Random-circuit control, 10 seeds; seed spread of our reservoir, 5 seeds & mean $\pm$ sd & R2\\
E5 & Shots $10^2,10^3,2{\times}10^3,10^4$ on Aer, with and without noise & MAE-vs-shots curve & R4\\
E6 & Hardware, ZF015 10--20 Aug, 88 circuits, Z basis, 10{,}000 shots (chosen from E5) & hardware vs noisy sim vs exact & R4\\
E7 & Ridge with and without the raw inputs added to the features & is the gain from the reservoir? & R3\\
E8 & ZF015 case study, 10--20 Aug 2023 & map: rebuilt vs predicted outline, time-to-reach town & R5\\
E9 & Sparse data: drop 15/30/45\% of hotspot days, carry last value & robustness curve, reservoir vs ESN & R1\\
E10 & Reconstruction check: rebuilt final outline vs official final perimeter, every 2023 fire & IoU; is the rebuild
sane? & R1\\
\bottomrule
\end{tabularx}
\end{center}

\begin{teamrule}
Pick the configuration in E3 on validation only. Freeze it. Then compute test numbers once. \supp{notes ch.~8}
\end{teamrule}

% =====================================================================
\section{Hardware plan}

\subsection{Budget}
\begin{center}\small
\begin{tabularx}{\linewidth}{@{}l Y@{}}
\toprule
\textbf{Item} & \textbf{Value}\\
\midrule
Plan limit & 10 min QPU per 28 days per account; job or batch mode, no \code{Session} \veri{IBM plans page; error 1352}\\
Circuits & ZF015, 10--20 Aug 2023: $11\times8=88$ (Z basis). Hotspots on all 11 days. \veri{re-run}\\
Shots & 2000 each $\Rightarrow176{,}000$ total\\
Shot noise & $\mathrm{SE}\le1/\sqrt{2000}=0.022$ per feature\\
Two-qubit gates & $L(n-1)=5\times4=20$ CZ, zero SWAPs\\
Rough fidelity & $(1-3.9{\times}10^{-3})^{20}\approx0.925$ (FakeFez median CZ error; two-qubit gates only) \veri{hardware pilot}
\caut{check ibm\_fez's calibration error}\\
\bottomrule
\end{tabularx}
\end{center}

\subsection{Steps}
\begin{enumerate}
\item Transpile one circuit; check CZ count $=20$. If higher, fix the layout before anything else.
\item \textbf{Pilot:} 4 circuits. Read QPU seconds used. Scale to the full circuit count. \veri{skipped: main job
  used 238\,s for 88 circuits $\times$ 10{,}000 shots (\code{job.usage()}, 3.95 resource units)}
\item If the estimate is $>5$ min: drop to 1000 shots, or keep only the 3 sectors facing the town.
\item \textbf{Main run:} all circuits in \textbf{one} \code{SamplerV2} job. Submitted by Pop, Thu 5 Nov 05:06\,UTC; done in under 5 min. \veri{job log}
\item Train ridge on Aer features (noise model of the same backend, same shots). Test on hardware features. \ours
\item If hardware error is much worse: run $\approx100$ training circuits on hardware and refit the ridge on hardware
  features. Report both.
\end{enumerate}

\begin{watch}
\begin{itemize}
\item \textbf{Shot noise is the biggest risk.} Stress test: NRMSE 0.18 exact $\to$ 1.01 at 100 shots.
  \supp{project3 \S9.1} Show the degradation; do not hide it.
\item \textbf{Fallback:} the main table is all simulation. The hardware row is extra, so a queue failure never breaks
  the demo. \supp{project3 \S9.3}
\item \textbf{Keys:} each member saves their own key with \code{getpass} in a terminal. Never in notebooks or git.
\end{itemize}
\end{watch}

% =====================================================================
\section{Team, interfaces, timeline}

\subsection{Roles}
\begin{center}\small
\begin{tabularx}{\linewidth}{@{}l l Y Y@{}}
\toprule
\textbf{ID} & \textbf{Role} & \textbf{Owns} & \textbf{Done when}\\
\midrule
R1 & Fire data & hotspots $\to$ fire groups $\to$ daily outlines $\to$ 8-sector radii and $\Delta d$; splits; E9, E10 &
\code{sectors.parquet} for 2018--2023\\
R2 & Quantum reservoir & chain port, exact + Aer features, E3, E4 & feature files for all splits\\
R3 & Weather + baselines + eval & per-fire weather from hotspots, $w_k$, scaling, baselines (incl.\ FBP), ridge,
metrics, E1, E7 & one script prints the main table\\
R4 & Hardware & noise model, E5, pilot, E6, budget log & hardware row + shots curve\\
R5 & Map + pitch & outline rebuild, time-to-reach, E2, E8 map, slides, Q\&A drill & 5 slides, rehearsed twice\\
\bottomrule
\end{tabularx}
\end{center}

\begin{watch}
R1 is the critical path: nothing runs until fire tracks exist. Build 2023 first (data already on disk: test + case
study), download 2018--2022 hotspots in parallel. R3 can build weather and the FBP baseline on 2023 at once.
\end{watch}

\subsection{Interfaces (agree tonight)}
Repo: \code{\textasciitilde/Projects/qff-ember/} with folders \code{data/ reservoir/ baselines/ eval/ hardware/
slides/}.
\begin{center}\small
\begin{tabularx}{\linewidth}{@{}l Y@{}}
\toprule
\textbf{File} & \textbf{Contents}\\
\midrule
\code{data/processed/hotspots\_<year>.parquet} & uid, rep\_date, lat, lon, ws, wd, rh, isi, bui, ros (points inside a
kept fire)\\
\code{data/processed/sectors.parquet} & uid, date, sector, r\_km, dd\_km, w\_k, isi, bui, rh, ros, n\_hot, masked, y1,
y3, split\\
\code{features/<model>\_<mode>.npz} & row ids, X; \ mode = exact, aer, aer\_noisy, ibm\\
\code{results/metrics.csv} & model, mode, lead, mae\_km, ss, iou, mae\_fast, n\_rows, n\_fires\\
\code{config.yaml} & $n$, $L$, $\gamma$, seed, leads, split years, input map, shots\\
\bottomrule
\end{tabularx}
\end{center}

\subsection{Timeline (Waterloo time)}
\begin{center}\small
\begin{tabularx}{\linewidth}{@{}l Y@{}}
\toprule
\textbf{When} & \textbf{Milestone}\\
\midrule
Wed eve & 30-min kick-off: read plan (verdict B), assign roles, agree interfaces.\\
Wed night & R1 2023 rebuilt outlines + sectors + E10 sanity check. R2 chain port passes exact-vs-NumPy test. R3 weather join + baselines on a
synthetic table. R4 noise model + pilot script. R5 deck skeleton + sourced case-study facts.\\
Thu 09:00 & R1 2018--2022 done. \textbf{Last switch point} to the FWI plan. R4 submits pilot, then the main job.\\
Thu 13:00 & \textbf{Checkpoint 1:} E1 table v1, E3 done, config frozen.\\
Thu 13--18 & E2, E4, E5, E7, E9; hardware results in; E8 map drafted.\\
Thu 18:00 & \textbf{Checkpoint 2:} every number in \code{metrics.csv}; figures drafted.\\
Thu 23:00 & \textbf{Code freeze.} Slides only after this.\\
Fri morning & Two timed rehearsals + Q\&A drill. \todo{R5: talk time and length from Discord}\\
\bottomrule
\end{tabularx}
\end{center}

Stretch, only if Checkpoint 2 is green: rank the towns around a fire by predicted time-to-reach (an
``evacuate first'' list).

% =====================================================================
\section{Expected outcomes}
Written \textbf{before} running, so the result cannot move the goalposts.

\begin{center}\small
\begin{tabularx}{\linewidth}{@{}Y Y@{}}
\toprule
\textbf{Expectation} & \textbf{Basis}\\
\midrule
Reservoir $\approx$ ESN, maybe slightly worse & project3: 0.181 vs 0.167 on NARMA-2 \supp{project3 \S9}\\
Persistence strong at 24\,h; weaker at 72\,h & fire speed changes with weather over days \ours\\
Growth comes in bursts (ZF015: 110 $\to$ 1{,}318 hotspots, 12 $\to$ 13 Aug) & hardest days to predict; check them
separately \veri{hotspot counts}\\
Any reservoir edge is at 24\,h, on the downwind side & reservoirs help at short range, fade by long range
\supp{QRC\_Broadcast\_Burns Fig.~1}\\
All models weak at 72\,h if the wind turns & no weather forecast in the inputs \ours\\
Hardware worse than noisy sim; refitting on hardware narrows the gap & size unknown \caut{report as found}\\
\bottomrule
\end{tabularx}
\end{center}

\textbf{Never say:} ``quantum advantage'', ``beats classical'', ``predicts wildfires'' (we predict how an existing fire
grows, not where fires start), ``replaces fire-spread models''.

% =====================================================================
\section{Results (test 2023, scored once)}
Configuration chosen on 2022 only (E3: $L=5$, $\gamma=\pi/4$, map $w_k$, ISI, $\Delta d_k$, BUI, RH), frozen with a hash
of \code{config.yaml}, then 2023 loaded and scored once. Features: exact noiseless expectation values. 392 fires.
\veri{metrics.csv, recomputed independently}

\begin{center}\small
\begin{tabular}{@{}l c c c c@{}}
\toprule
& \multicolumn{2}{c}{\textbf{+24\,h}} & \multicolumn{2}{c}{\textbf{+72\,h}}\\
\textbf{Model} & MAE (km) [95\% CI] & SS & MAE (km) [95\% CI] & SS\\
\midrule
Persistence & 0.349 [0.307, 0.389] & 0 & 1.070 [0.933, 1.196] & 0\\
FBP physics & \textbf{0.262} [0.233, 0.289] & \textbf{0.25} & \textbf{0.694} [0.609, 0.772] & \textbf{0.35}\\
Raw window ridge & 0.348 [0.319, 0.375] & 0.00 & 0.920 [0.840, 0.990] & 0.14\\
ESN, 20 units & 0.350 [0.322, 0.376] & 0.00 & 0.934 [0.858, 1.004] & 0.13\\
\textbf{Quantum reservoir} & 0.360 [0.330, 0.387] & $-0.03$ & 0.943 [0.864, 1.015] & 0.12\\
Reservoir + raw (E7) & 0.348 [0.318, 0.374] & 0.00 & 0.921 [0.843, 0.991] & 0.14\\
Random circuit, 10 seeds & 0.361 & $-0.03$ & 0.946 & 0.12\\
\bottomrule
\end{tabular}
\end{center}

Shape and growing sectors (IoU of the predicted vs true octagon, higher is better; MAE on sectors with $y>0$):
\begin{center}\small
\begin{tabular}{@{}l c c c c@{}}
\toprule
& \multicolumn{2}{c}{\textbf{IoU}} & \multicolumn{2}{c}{\textbf{MAE, growing (km)}}\\
\textbf{Model} & +24\,h & +72\,h & +24\,h & +72\,h\\
\midrule
Persistence & \textbf{0.918} & \textbf{0.832} & 1.705 & 3.637\\
FBP physics & 0.781 & 0.710 & 1.593 & 2.690\\
Raw window ridge & 0.766 & 0.672 & 1.425 & 2.365\\
ESN, 20 units & 0.763 & 0.669 & 1.429 & 2.364\\
\textbf{Quantum reservoir} & 0.757 & 0.663 & 1.456 & 2.389\\
Reservoir + raw (E7) & 0.765 & 0.671 & \textbf{1.421} & \textbf{2.360}\\
Random circuit & 0.761 & 0.666 & 1.446 & 2.381\\
\bottomrule
\end{tabular}
\end{center}

Paired fire-bootstrap differences, reservoir minus reference (negative = reservoir better):
\begin{center}\small
\begin{tabular}{@{}l c c@{}}
\toprule
\textbf{Reference} & \textbf{+24\,h} & \textbf{+72\,h}\\
\midrule
Persistence & $+0.011$ [$-0.004$, 0.025] & $\mathbf{-0.127}$ [$-0.186$, $-0.067$]\\
Raw window ridge & $+0.012$ [0.010, 0.014] & $+0.023$ [0.017, 0.029]\\
ESN, 20 units & $+0.010$ [0.006, 0.013] & $+0.008$ [$-0.002$, 0.018]\\
\bottomrule
\end{tabular}
\end{center}

\begin{itemize}
\item FBP has the lowest MAE at both leads. It uses the modelled spread rate \code{ros}; the reservoir inputs do not.
\item No model wins on every measure (table below). Persistence has the best outline overlap: most sectors do not
  move, and it predicts exactly 0 there. On growing sectors only, every learned model beats FBP and persistence.
\item The reservoir beats persistence only at +72\,h, and so does every ridge model. It trails the raw ridge at both
  leads; the reservoir adds nothing over its own inputs (E7).
\item Seeds (E4): reservoir $0.362\pm0.008$ / $0.950\pm0.016$, random circuit $0.361\pm0.007$ / $0.946\pm0.016$.
  The chain circuit is no better than a random one.
\item E2: at +24\,h the reservoir loses to persistence downwind, crosswind and upwind; at +72\,h it gains in all three.
\item Decision (Pop, 5 Nov): report as is. No hybrid model, no second look at 2023.
\end{itemize}

\subsection{E6: real hardware (ibm\_fez, ZF015, 10--20 Aug 2023)}
Job \code{dbm0tm58ft8s738tu920}: 88 circuits, 10{,}000 shots each, 20 CZ, no SWAP. QPU time 238\,s (of 600\,s).
Frozen exact-trained Z-only read-out, applied unchanged to each feature source. 88 sector windows from one fire.
\begin{center}\small
\begin{tabular}{@{}l c c c@{}}
\toprule
\textbf{Feature source} & \textbf{Feature error} (mean $|\Delta|$) & \textbf{+24\,h MAE (km)} & \textbf{+72\,h MAE (km)}\\
\midrule
Exact (simulator) & 0 & 0.272 & 1.001\\
FakeFez (noisy simulator) & 0.044 & 0.277 & 1.006\\
\textbf{ibm\_fez hardware} & 0.073 & 0.258 & 0.979\\
Persistence & -- & 0.301 & 1.228\\
\bottomrule
\end{tabular}
\end{center}
\begin{itemize}
\item The hardware features follow the exact ones closely, but they shrink towards 0 at large $|\langle Z\rangle|$
  (decay and readout error).
\item The forecasts survive the noise: hardware MAE is within 0.02\,km of exact at both leads.
\item The hardware being slightly lower is not a real gain: it is one fire, the 88 rows are correlated, and 93\% of
  the +24\,h targets are 0. No CI is given and nothing is claimed beyond this case.
\end{itemize}

% =====================================================================
\section{Pitch and questions}

\subsection{Five slides}
\begin{center}\small
\begin{tabularx}{\linewidth}{@{}l l Y@{}}
\toprule
\textbf{\#} & \textbf{Slide} & \textbf{Content}\\
\midrule
1 & Problem & 2023 numbers, Yellowknife evacuation, ``where does it go next?''\\
2 & Task & outlines rebuilt from satellite hotspots; 8 sectors, advance at 24\,h and 72\,h, from day 3\\
3 & Method & circuit diagram, inputs per qubit, only read-out trained, ``extends PronoiaQ''\\
4 & Results & main table, downwind vs upwind plot, 2023 map with predicted outline\\
5 & Hardware + limits & hardware row, shots curve, QPU seconds; parity not advantage; next steps\\
\bottomrule
\end{tabularx}
\end{center}

\subsection{Questions to rehearse}
\begin{center}\small
\begin{tabularx}{\linewidth}{@{}Y Y@{}}
\toprule
\textbf{Question} & \textbf{Answer from}\\
\midrule
Where is the nonlinearity, if quantum dynamics is linear? & encoding + products of inputs (notes ch.~3)\\
How does it predict day 2 and day 3 if qubits cannot be reset? & new run per sample; separate ridge per lead
(Section 4)\\
Why not just an ESN? & E1 table, same read-out size, fully connected\\
Why not the operational spread models? & Section 5, last box\\
Shots, two-qubit gates, QPU seconds? & Section 7 budget table\\
Why 8 sectors and not a map? & Section 3 decision box\\
Your target comes from the same hotspots as your input? & yes; Section 3 limitation table, E10\\
What is new this week; who built what? & Section 2 table + roles (Section 8)\\
\bottomrule
\end{tabularx}
\end{center}

% =====================================================================
\section{Soundness audit}

\begin{center}\small
\begin{tabularx}{\linewidth}{@{}Y l Y@{}}
\toprule
\textbf{Claim} & \textbf{Status} & \textbf{Basis}\\
\midrule
Perimeter archives are daily & \textbf{no} & perimeter-check.md, re-run (4 Nov)\\
Hotspots are dated and cover ZF015 daily in 10--20 Aug & supported & re-run counts (4 Nov)\\
Rebuilt outlines track the real edge & partly & E10: median IoU 0.65 / 0.83 buffered\\
Hotspot weather represents the fire's weather & supported (modelled) & hotspot attribute sheet\\
Chain needs no SWAPs on heavy-hex & \textbf{yes} & FakeFez transpile: CZ $=20$, depth 41\\
Reservoir matches ESN here & \textbf{close, slightly worse} & E1: $+0.010$\,km at +24\,h (paired CI 0.006--0.013)\\
Reservoir helps at 24\,h downwind & \textbf{no} & E1, E2: loses to persistence in every wind bin at +24\,h\\
Hardware reproduces simulation & yes, on one fire & E6: MAE within 0.02\,km of exact\\
Quantum advantage & \textbf{not claimed} & no source supports it\\
\bottomrule
\end{tabularx}
\end{center}

% =====================================================================
\phantomsection\hypertarget{refs}{}
\section*{References}
\markboth{References}{}

\subsection*{Data}
\begin{enumerate}[label={[D\arabic*]},leftmargin=3em]
\item Natural Resources Canada, Canadian Forest Service. \emph{Canadian Wildland Fire Information System (CWFIS):
  station Fire Weather Index observations}, files \code{cwfis\_fwi2000sv3.0\_ll.csv.bz2},
  \code{cwfis\_fwi2010sv3.0\_ll.csv.bz2}, \code{cwfis\_fwi2020sv3.0\_ll.csv.bz2}.
  \url{https://cwfis.cfs.nrcan.gc.ca/downloads/fwi_obs/} (accessed 4 Nov 2026).
\item Natural Resources Canada, Canadian Forest Service. \emph{CWFIS reporting weather stations list},
  \code{cwfis\_allstn2025.csv}, same directory (accessed 4 Nov 2026). Catalogue record
  71e96091-9d23-4963-982d-9df7c7e0ea90:
  \url{https://cwfis.cfs.nrcan.gc.ca/en/catalogue/results/71e96091-9d23-4963-982d-9df7c7e0ea90}.
\item Natural Resources Canada, Canadian Forest Service. \emph{CWFIS fire perimeter and hotspot archives},
  \code{<YEAR>\_perimeters.zip} (2018--2023), \code{<YEAR>\_hotspots.zip} (2012--2023).
  \url{https://cwfis.cfs.nrcan.gc.ca/downloads/hotspots/archive/} (accessed 4 Nov 2026). Attribute sheets:
  \url{https://cwfis.cfs.nrcan.gc.ca/downloads/hotspots/fireM3_hotspots_metadata_attributes_EN_FR.pdf},
  \url{https://cwfis.cfs.nrcan.gc.ca/downloads/hotspots/fireM3_buffered_hotspots_metadata_attributes_EN_FR.pdf}.
\item Government of the Northwest Territories, \emph{Fire History} map layer 81 (fire IDs ZF015-23 etc.),
  \url{https://www.apps.geomatics.gov.nt.ca/arcgis/rest/services/GNWT/BiologicEcologic_LCC/MapServer/81}
  (accessed 4 Nov 2026). Used to match CWFIS uid 242981 to ZF015-23.
\end{enumerate}

\subsection*{Fire behaviour}
\begin{enumerate}[label={[F\arabic*]},leftmargin=3em]
\item C.~E. Van Wagner, \emph{Development and Structure of the Canadian Forest Fire Weather Index System},
  Forestry Technical Report 35, Canadian Forestry Service, Ottawa (1987). Defines ISI and BUI.
\item Forestry Canada Fire Danger Group, \emph{Development and Structure of the Canadian Forest Fire Behavior
  Prediction System}, Information Report ST-X-3, Ottawa (1992). The operational spread model we do not replace.
  \caut{check the citation before using on a slide}
\end{enumerate}

\subsection*{Method}
\begin{enumerate}[label={[M\arabic*]},leftmargin=3em]
\item K.~Fujii, K.~Nakajima, ``Harnessing disordered-ensemble quantum dynamics for machine learning,''
  \emph{Phys.\ Rev.\ Applied} \textbf{8}, 024030 (2017).
\item P.~Mujal \emph{et al.}, ``Opportunities in quantum reservoir computing and extreme learning machines,''
  \emph{Adv.\ Quantum Technol.} \textbf{4}, 2100027 (2021).
\item A.~Sannia \emph{et al.}, ``Dissipation as a resource for quantum reservoir computing,''
  \emph{Quantum} \textbf{8}, 1291 (2024).
\item H.~Jaeger, H.~Haas, ``Harnessing nonlinearity: predicting chaotic systems and saving energy in wireless
  communication,'' \emph{Science} \textbf{304}, 78--80 (2004). (ESN baseline.)
\end{enumerate}

\subsection*{Related and own prior work}
\begin{enumerate}[label={[P\arabic*]},leftmargin=3em]
\item PronoiaQ team, Quantathon 2026 code and \code{CONCLUSION.md} (windowed restart reservoir, Qiskit twin,
  multi-lead read-out; marine-heatwave results).
\item \emph{Quantum Reservoir Computing for Forecasting Nonlinear Clean-Energy and Climate Time Series},
  hackathon specification with stress test (30 May 2026). Source of the ESN-parity and shot-noise numbers.
\item GitHub user Luke-Fu101, \code{QRC\_Broadcast\_Burns}, \url{https://github.com/Luke-Fu101/QRC_Broadcast_Burns}
  (accessed 4 Nov 2026). Reservoir vs ESN on Lorenz-63 for prescribed-burn weather windows; simulation only.
\end{enumerate}

\end{document}
